\documentclass[../NDCNNAI_main.tex]{subfiles}
\graphicspath{{\subfix{../fig/}}}

\begin{document}

\subsection*{Функциональные пространства, операторы и метрики}

Пусть $D \subset \mathbb{R}^d$, $U \subset \mathbb{R}^n$ --- компактные (то есть замкнутые и ограниченные) множества.

\begin{definition}[Пространства $C(D)$, $C(U)$ и супремум-норма]
$C(D)$ и $C(U)$ --- пространства вещественных непрерывных функций на $D$ и $U$.
Супремум-норма задаётся формулами
$$
\|f\|_\infty \coloneq \sup_{x\in D} |f(x)|,\qquad
\|g\|_\infty \coloneq \sup_{y\in U} |g(y)|.
$$
\end{definition}

\begin{remark}
Эта норма индуцирует метрику $d(f,g)=\|f-g\|_\infty$, поэтому $C(D)$ и $C(U)$ являются банаховыми пространствами.
\end{remark}

\begin{definition}[Пространства $L^p(D)$ и $L^p$-норма]
Для $1 \le p < \infty$ пространство $L^p(D)$ состоит из (классов эквивалентности) измеримых функций с нормой
$$
\|f\|_{L^p(D)} \coloneq \left(\int_D |f(x)|^p\,dx\right)^{1/p}.
$$
\end{definition}

\begin{definition}[Скалярное произведение в $L^2(U)$]
Пространство $L^2(U)$ является гильбертовым со скалярным произведением
$$
\langle f,g\rangle_{L^2(U)} \coloneq \int_U f(y)\,g(y)\,dy,
$$
и соответствующей нормой
$$
\|g\|_{L^2(U)} = \langle g,g\rangle_{L^2(U)}^{1/2}.
$$
\end{definition}

\begin{definition}[Оператор между пространствами функций]
Оператор --- это отображение $G\colon X \to Y$ между пространствами функций (банаховыми или гильбертовыми), например
$$
G\colon C(D) \to C(U)
\quad \text{или} \quad
G\colon C(D) \to L^2(U).
$$
\end{definition}

Будем говорить о $X$ как о пространстве входных функций $u$, а о $Y$ --- как о пространстве выходных функций $G(u)$.
DeepONets будем рассматривать как параметризованное семейство таких операторов.

\begin{definition}[Вероятностная мера на $X$]
Вероятностная мера $\mu$ на $X$ задаёт распределение случайного элемента
$u\colon\Omega \to X$ (то есть закон случайной функции $u \in X$, это можно интерпретировать как распределение обучающих данных.)
\end{definition}
\begin{definition}[Прямой образ меры $G\#\mu$] Прямой образ (push-forward) $G\#\mu$ --- это распределение случайного элемента $G(u)$ при $u \sim \mu$:
$$
(G\#\mu)(B) \coloneq \mu\!\left(G^{-1}(B)\right),
\qquad B \subset Y \ \text{--- борелевское множество}.
$$
\end{definition}

\begin{definition}[Нецентрированный оператор ковариации]
Пусть $(H,\langle\cdot,\cdot\rangle_H)$ --- сепарабельное гильбертово пространство и
$\nu \in \mathcal{P}(H)$ такова, что
$$
\int_H \|z\|_H^2\, d\nu(z) < \infty.
$$
Тогда (нецентрированный) оператор ковариации $C_\nu\colon H \to H$ определяется как
$$
C_\nu v \coloneq \int_H \langle z, v\rangle_H\, z\, d\nu(z).
$$
\end{definition}

\begin{remark}
Оператор $C_\nu$ является самосопряжённым, положительным и оператором следового класса; позже мы будем использовать его собственные пары (собственное значение–собственный вектор) для построения хороших базисов trunk-ветви DeepONets.
\end{remark}
% Напомним определение оператора следового класса.
% \begin{definition}[Оператор следового класса]
% Пусть $H$ --- (сепарабельное) гильбертово пространство, а $A\colon H\to H$ --- ограниченный линейный оператор.
% Обозначим $|A| \coloneq (A^*A)^{1/2}$ и пусть $\{s_k(A)\}_{k\ge 1}$ --- сингулярные числа оператора $A$
% (то есть собственные значения $|A|$, упорядоченные по убыванию с учётом кратности).
% Оператор $A$ называется \emph{оператором следового класса} (trace-class), если
% $$
% \sum_{k=1}^{\infty} s_k(A) < \infty.
% $$
% В этом случае определена \emph{следовая норма}
% $$
% \|A\|_1 \coloneq \mathrm{Tr}(|A|) = \sum_{k=1}^{\infty} s_k(A) < \infty.
% $$
% \end{definition}

% \begin{remark}
% Если $A$ --- самосопряжённый неотрицательный оператор, то
% $$
% \mathrm{Tr}(A)=\sum_{k=1}^{\infty} \lambda_k(A),
% $$
% где $\lambda_k(A)$ --- его собственные значения (с учётом кратности).
% Эквивалентно, для любого ортонормированного базиса $\{e_k\}$ выполнено
% $$
% \mathrm{Tr}(A)=\sum_{k=1}^{\infty} \langle A e_k, e_k\rangle_H,
% $$
% и эта сумма не зависит от выбора базиса.

% \end{remark}

\subsection*{Обозначения для нейронных сетей (конечномерные отображения)}

Мы будем неоднократно использовать стандартные полносвязные нейронные сети прямого распространения (feedforward) как базовые (конечномерные) блоки.

\begin{definition}[Полносвязная feedforward нейронная сеть]
Feedforward нейронная сеть $L\colon \mathbb{R}^{d_{\mathrm{in}}}\to \mathbb{R}^{d_{\mathrm{out}}}$ с $K-1$ скрытыми слоями задаётся соотношениями
$$
z_1 = x,\qquad
z_{k+1} = \sigma(W_k z_k + b_k),\quad k=1,\ldots,K-1,\qquad
L(x)= W_K z_K + b_K,
$$
где
$$
W_k \in \mathbb{R}^{m_k\times m_{k-1}},\qquad
b_k \in \mathbb{R}^{m_k},\qquad
m_0=d_{\mathrm{in}},\quad m_K=d_{\mathrm{out}}.
$$
\end{definition}

\textbf{Параметры архитектуры.}
\begin{itemize}
  \item \emph{Ширина} $k$-го слоя: $m_k$ (число нейронов в этом слое).
  \item \emph{Глубина}: число скрытых слоёв $=K-1$.
  \item \emph{Параметры}: $\theta=\{W_k,b_k\}_{k=1}^{K}$.
  \item \emph{Размер} (общее число скалярных параметров):
  $$
  \mathrm{size}(L)=\|\theta\|_{\ell^0}
  =\sum_{k=1}^{K}\bigl(m_k m_{k-1}+m_k\bigr).
  $$
\end{itemize}

\subsection{Асимптотические обозначения: $O(\cdot)$ и $\Theta(\cdot)$}

Мы часто будем описывать ``размер'' сети (число слоёв/нейронов) с точностью до постоянных множителей (констант), используя асимптотические обозначения.

Пусть далее $f,g\colon \mathbb{N}\to (0,\infty)$.

\begin{definition}[$O$-большое, верхняя оценка]
Говорят, что $f(k)=O(g(k))$, если существуют константы $C>0$ и $k_0\in\mathbb{N}$ такие, что
$$
f(k)\le C\,g(k)\qquad \text{для всех } k\ge k_0.
$$
\end{definition}

\begin{remark}
Возможна следующая интерпретация: при больших $k$ функция $f$ растёт не быстрее, чем $g$, с точностью до константы.
\end{remark}

\begin{definition}[$\Theta$-большое, точная (двусторонняя) оценка]
Говорят, что $f(k)=\Theta(g(k))$, если существуют константы $c_1,c_2>0$ и $k_0\in\mathbb{N}$ такие, что
$$
c_1\,g(k)\le f(k)\le c_2\,g(k)\qquad \text{для всех } k\ge k_0.
$$
\end{definition}

\begin{remark}
Можно интерпретировать так: при больших $k$ функция $f$ растёт как $g$ с точностью до константы
(то есть выполняются и верхняя, и нижняя оценки).
\end{remark}

\begin{example}
  Приведём некоторые примеры использования асимптотической нотации.
  \begin{itemize}
  \item $3k^2+5k+7=\Theta(k^2)$ и, следовательно, также $O(k^2)$.
  \item $2^k$ не является $O(k^m)$ ни для какой фиксированной степени $m$:
  экспоненциальная функция растёт быстрее любой степенной (полиномиальной).
  \item ``$\Theta(k^3)$ слоёв'' означает, что число слоёв лежит между $c_1k^3$ и $c_2k^3$ для некоторых констант $c_1,c_2>0$.
  \item ``ширина $O(1)$'' означает, что ширина ограничена константой $C$, не зависящей от $k$.
\end{itemize}
\end{example}

\subsection{Разделение по глубине (Тельгарский)}

\begin{theorem}[Об иерархии по глубине]
Для каждого $k\in\mathbb{N}$ существует сеть с ReLU-активацией (ReLU-сеть) $f_k\colon [0,1]^d \to \mathbb{R}$ с $\Theta(k^3)$ слоями и шириной $O(1)$ такая, что любая ReLU-сеть с не более чем $k$ скрытыми слоями, которая аппроксимирует $f_k$ в $L^1([0,1]^d)$ с ошибкой $<\frac{1}{8}$, должна иметь по крайней мере $\Omega(2^k)$ нейронов.
\end{theorem}

\begin{proof}[Идея построения (одномерный случай)]
\leavevmode
\begin{enumerate}
  \item Определим ``треугольную'' функцию на $[0,1]$:
  $$
  \Delta(x)=\max\{0,\, 1-|2x-1|\},
  $$
  которую можно реализовать очень небольшой ReLU-сетью.
  \item Применим её к самой себе $k$ раз:
  $$
  x \mapsto \Delta^{k}(x)=\Delta(\Delta(\ldots \Delta(x)\ldots)).
  $$
  Каждая композиция удваивает число пиков, то есть получается порядка $2^k$ осцилляций (много чередующихся линейных участков).
  \item Расширим до $[0,1]^d$, игнорируя лишние координаты:
  $$
  f_k(x_1,\ldots,x_d)=\Delta^{k}(x_1),
  $$
  при этом ширина остаётся $O(1)$.
\end{enumerate}
\textbf{Почему не могут аппроксимировать $\Delta^{k}$ без экспоненциального роста ширины?}
\begin{itemize}
  \item ReLU-сеть глубины $k$ с $W$ нейронами задаёт кусочно-аффинное отображение:
  пространство входов разбивается на области, на каждой из которых сеть совпадает с некоторой аффинной функцией
  (то есть имеет вид $x \mapsto Ax+b$).
  Число таких областей в общем случае не превосходит $O(W^k)$.
  \item Функция $\Delta^{k}$ создаёт порядка $2^k$ чередующихся ``участков поведения''
  (интуитивно: порядка $2^k$ пиков/колебаний, или много попеременно возрастающих и убывающих линейных отрезков).
1  Чтобы аппроксимировать такую структуру неглубокой ReLU-сетью, ей нужно иметь достаточно много аффинных областей,
  а значит требуется как минимум $W \ge c\,2^k$ нейронов.
\end{itemize}

\end{proof}

\begin{remark}[Неформальный вывод]
Глубина может давать экспоненциальную экономию числа нейронов (и параметров) при представлении некоторых функций:
глубокие сети способны эффективно задавать сильно осциллирующие функции (быстро меняющиеся, с большим числом колебаний),
тогда как неглубоким сетям для этого требуется экспоненциально больше нейронов.

Это объясняет, почему в архитектурах DeepONets имеет смысл использовать достаточно глубокие сети в ветвях branch и trunk:
глубина может заметно повысить эффективность представления сложных зависимостей без чрезмерного роста числа параметров.
\end{remark}

\subsection{Достаточность ограниченной ширины (Ханин)}

\begin{theorem}[УА при ограниченной ширине (неформально)]
Пусть $d_{\mathrm{in}}\ge 1$ и $K\subset \mathbb{R}^{d_{\mathrm{in}}}$ --- компакт. Тогда:
\begin{itemize}
  \item Если ReLU-сеть имеет ширину не меньше $d_{\mathrm{in}}+1$, то для любой непрерывной функции
  $f\colon K\to \mathbb{R}$ и любого $\varepsilon>0$ существует ReLU-сеть $L$ ширины $d_{\mathrm{in}}+1$
  (глубина не ограничена), такая что
  $$
  \sup_{x\in K} |f(x)-L(x)|<\varepsilon.
  $$
  \item Если ширина $\le d_{\mathrm{in}}$, то существует непрерывная функция $f\colon K\to\mathbb{R}$ и число
  $\varepsilon_0>0$ такие, что для любой ReLU-сети $L$ этой ширины (при любой глубине) выполнено
  $$
  \sup_{x\in K}|f(x)-L(x)|\ge \varepsilon_0.
  $$
  (более неформально : если ширина $\le d_{\mathrm{in}}$, то как бы мы ни увеличивали глубину сети,
  не удаётся аппроксимировать все непрерывные функции на $K$.)
\end{itemize}
\end{theorem}


Приведём некоторые \textbf{пояснение} к этому утверждению.
\begin{itemize}
  \item Ширина $d_{\mathrm{in}}+1$ (а в ряде формулировок теоремы --- $d_{\mathrm{in}}+2$) по существу является минимальной,
  при которой ReLU-сети обладают свойством универсальной аппроксимации:
  они могут равномерно приближать любую непрерывную функцию $f\colon K\to\mathbb{R}$ с любой заданной точностью.
  \item После того как ширина превышает этот порог, точность аппроксимации можно улучшать, увеличивая глубину сети,
  не увеличивая её ширину.
\end{itemize}

\textbf{Геометрическая интуиция.}
\begin{itemize}
  \item ReLU-сети задают кусочно-аффинные отображения: существует конечное разбиение множества $K$ на области
  $K_1,\dots,K_M$ (пересечения $K$ с многогранниками) такое, что на каждой области
  $$
  L(x)=A_jx+b_j,\qquad x\in K_j.
  $$
  \item При ширине $\le d_{\mathrm{in}}$ возникает препятствие:
  класс функций, реализуемых ReLU-сетями такой ширины (при любой глубине), не является плотным в $C(K)$ по норме
  $\|\cdot\|_\infty$, поэтому не все непрерывные функции на $K$ можно аппроксимировать.
  \item Если увеличить ширину до $d_{\mathrm{in}}+1$ (то есть добавить один нейрон),
  это препятствие исчезает: увеличивая глубину, можно строить сколь угодно \emph{мелкие} разбиения множества $K$
(то есть разбиения на части малого размера), что и позволяет приближать произвольные
$f\in C(K)$.

\end{itemize}

\begin{remark}[Практический вывод]
В задачах обучения операторов можно держать сети branch и trunk \emph{узкими}
(ширины порядка $d_{\mathrm{in}}+1$) и увеличивать глубину, чтобы описывать сложные зависимости.
Это особенно удобно, когда входная размерность велика; в задачах обучения операторов это часто означает,
что вход задаётся большим числом измерений  $u$ в различных точках.
\end{remark}

\subsection{От операторных сетей к DeepONets}

\textbf{Классические операторные сети (Чен--Чен).}
Пусть $G\colon K \subset C(D)\to C(U)$ непрерывен на компакте $K$.
Чен--Чен аппроксимируют $G$ конечной суммой слагаемых разделённого вида:
$$
G(u)(y)\approx \sum_{k=1}^{p}\beta_k(u)\,\tau_k(y),
$$
то есть суммой членов, где множитель $\beta_k(u)$ зависит только от $u$, а множитель $\tau_k(y)$ --- только от $y$.
Здесь
$$
\beta_k(u)=\sum_{i=1}^{n} c_i^{k}\,\sigma\!\left(\sum_{j=1}^{m}\xi_{ij}^{k}\,u(x_j)+\theta_i^{k}\right),
\qquad
\tau_k(y)=\sigma\!\left(w_k\cdot y+\zeta_k\right).
$$
Ключевые моменты такого подхода:
\begin{itemize}
  \item Входная функция $u\in C(D)$ сначала берётся в точках $x_1,\ldots,x_m\in D$: в сеть поступает только вектор значений
  $(u(x_1),\ldots,u(x_m))$.
  \item Для каждого $k$ коэффициент $\beta_k(u)$ вычисляется неглубокой полносвязной сетью
  (MLP, \emph{многослойным перцептроном}) по входному вектору $(u(x_1),\ldots,u(x_m))$.
  \item Функции $\tau_k\colon U\to\mathbb{R}$ также задаются неглубокими сетями,
  принимающими на вход пространственную переменную $y$.
  \item Поэтому аппроксимация имеет \emph{низкоранговый} вид: это сумма $p$ разделимых произведений
  $\beta_k(u)\tau_k(y)$.
  В терминологии DeepONets функции $\beta_k(u)$ соответствуют \emph{branch}-части,
  а функции $\tau_k(y)$ --- \emph{trunk}-части.
\end{itemize}


\textbf{Точка зрения DeepONet.}
DeepONets сохраняют ту же общую схему разложения, что и у Чен--Чен,
но используют более глубокие нейросети и явно выделяют архитектуру как композицию модулей:
\begin{align*}
  & E\colon C(D)\to \mathbb{R}^m &  \text{(энкодер; вход branch-ветви)},\\
  & A\colon \mathbb{R}^m \to \mathbb{R}^p & \text{(сеть, вычисляющая коэффициенты)},\\
  & R\colon \mathbb{R}^p \to C(U) & \text{(декодер; trunk-ветвь)},
\end{align*}
так что
$$
N(u) = (R\circ A\circ E)(u).
$$

\begin{itemize}
  \item $E$ переводит входную функцию $u$ в конечномерный вектор признаков.
  Обычно $E(u)$ состоит из точечных значений $u(x_1),\ldots,u(x_m)$, но в более общем случае
  $E$ может включать и другие линейные функционалы от $u$ (например, интегралы или свёртки).
  \item $A$ --- глубокий MLP в $\mathbb{R}^m$, который по набору признаков $E(u)$ выдаёт вектор коэффициентов
  $(\beta_1(u),\ldots,\beta_p(u))$.
  \item $R$ по коэффициентам восстанавливает функцию на $U$; например,
  $$
  R(\alpha)(y)=\tau_0(y)+\sum_{k=1}^{p}\alpha_k\,\tau_k(y),
  $$
  где функции $\tau_k$ параметризуются trunk-сетью, то есть вычисляются как выходы нейросети от переменной $y$.
\end{itemize}

\begin{remark}[Ключевая идея]
DeepONet $N=R\circ A\circ E$ можно рассматривать как глубокую версию операторной сети Chen--Chen:
сохраняется разложение ``коэффициенты от $u$ + функции от $y$'', но используются более выразительные
(глубокие) branch- и trunk-сети, а вся конструкция имеет прозрачную модульную структуру
$E \Rightarrow A \Rightarrow R$.
\end{remark}
Наглядно архитектура DeepONet, записанная как композиция $N=R\circ A\circ E$, изображена на рис.~\ref{fig:deeponet}.

\begin{figure}[ht!]
\centering
\begin{adjustbox}{max width=\textwidth}
\begin{tikzpicture}[
  >=Stealth,
  font=\small,
  node distance=10mm and 12mm,
  box/.style={draw, rounded corners, align=center, inner sep=5pt, minimum height=12mm},
  lab/.style={align=center}
]
\node[lab] (uin) {$u\in C(D)$};

\node[box, right=16mm of uin, minimum width=42mm] (enc)
{Encoder\\[1mm] $E:\ u\mapsto \bigl(u(x_1),\ldots,u(x_m)\bigr)$};

\node[box, right=12mm of enc, minimum width=32mm] (app)
{Approximator\\[1mm] $A:\ \mathbb{R}^m\to \mathbb{R}^p$};

\node[box, below=14mm of app, minimum width=74mm] (rec)
{Reconstructor\\[1mm]
$R(\alpha)(y)=\tau_0(y)+\displaystyle\sum_{k=1}^{p}\alpha_k\,\tau_k(y)$};

\node[lab, right=14mm of rec] (nout) {$N(u)\in C(U)$};

\draw[->] (uin) -- (enc);
\draw[->] (enc) -- node[above] {$\mathbb{R}^m$} (app);
\draw[->] (app) -- node[right] {$\mathbb{R}^p$} (rec);
\draw[->] (rec) -- (nout);
\end{tikzpicture}
\end{adjustbox}
\caption{Схема DeepONet: $N=R\circ A\circ E$.}
\label{fig:deeponet}
\end{figure}

Вспомним теперь результат, доказанный в \hyperref[thm:chen_chen_oper]{лекции 3}, он пригодится для вывода похожего утверждения для DeepONet.
\begin{theorem}[УАТ для операторов, Чен--Чен]
Пусть $X$ --- банахово пространство, $K_1\subset X$, $K_2\subset \mathbb{R}^d$ --- компакты,
и $V\subset C(K_1)$ --- компактное множество входных функций.
Предположим, что оператор $G\colon V\to C(K_2)$ непрерывен, а функция активации
$\sigma\colon \mathbb{R}\to\mathbb{R}$ непрерывна и не является многочленом.
Тогда для любого $\varepsilon>0$ существуют целые числа $m,p$,
$x_1,\ldots,x_m\in K_1$ и сети $\{\beta_k\}_{k=1}^{p}$, $\{\tau_k\}_{k=1}^{p}$ формы Чен--Чена такие, что
$$
\sup_{u\in V,\ y\in K_2}\left|\,G(u)(y)-\sum_{k=1}^{p}\beta_k(u)\,\tau_k(y)\right|<\varepsilon.
$$
\end{theorem}

Краткое \textbf{пояснение.}
\begin{itemize}
  \item Мы аппроксимируем оператор $u\mapsto G(u)$ равномерно по всем $u\in V$ и $y\in K_2$.
  \item Аппроксимация имеет разделимый (branch--trunk) вид $\sum_k \beta_k(u)\tau_k(y)$,
  причём обе части реализуются стандартными (неглубокими) нейронными сетями.
  \item Это операторный аналог классической теоремы универсальной аппроксимации на $C(K)$.
\end{itemize}
Напомним \textbf{идею доказательства (три шага).}
\begin{enumerate}
  \item Дискретизуем входную функцию $u$ в конечном числе точек $x_j$ и рассматриваем вектор
  \[
    E(u)\coloneq (u(x_1),\ldots,u(x_m))\in\mathbb{R}^m.
  \]
  На компакте $V$ отображение $E\colon V\to\mathbb{R}^m$ непрерывно.

  \item Для конечномерного отображения $E(u)\mapsto (\beta_1(u),\ldots,\beta_p(u))$ применяем стандартную теорему
  об универсальной аппроксимации на компактах в $\mathbb{R}^m$ (Цыбенко/Хорник), получая аппроксимацию
  с помощью (неглубокой) полносвязной сети.

  \item Зависимость по $y\in K_2$ аппроксимируем общей для всех $u$ системой функций $\{\tau_k\}$,
  а затем ``склеиваем'' аппроксимации, используя равномерность оценок на компактах.
\end{enumerate}

\end{document}
